\documentclass[10pt,a4paper]{amsart}
\usepackage[margin=19mm]{geometry}
\usepackage{amsmath,amssymb,mathtools}
\usepackage[scr=rsfs,cal=euler]{mathalfa}
\usepackage{xfrac}
\usepackage[unicode,hidelinks,bookmarks=false]{hyperref}
\newcommand{\ie}{i.e.\ }
\newcommand{\cf}{cf.\ }
\newcommand{\resp}{respectively\ }
\newcommand{\Subsection}{\subsection}
\providecommand{\nobreakdash}{\nobreak\hbox{-}}
\setlength{\emergencystretch}{2em}
\multlinegap0pt
\usepackage{bm}
\usepackage{bbm}
\usepackage{dsfont}
\usepackage{xspace}
\usepackage{cancel}
\usepackage{mathtools}
\usepackage{amsthm}
\makeatletter
\@ifundefined{Prop}{\newtheorem{Prop}{Proposition}}{}
\@ifundefined{lemma}{\newtheorem{lemma}{Lemma}}{}
\makeatother
\newcommand{\Ord}{\mathop {\rm \bf ord}}
\title[Normal forms for ordinary differential operators, I]{Normal forms for ordinary differential operators, I}
\author{J. Guo, A.B. Zheglov}
\date{Source: arXiv:2406.14414v4 (CC BY 4.0)}
\begin{document}
\maketitle

\noindent\emph{Source and licence.} Normal forms for ordinary differential operators, I. Version arXiv:2406.14414v4 (CC BY 4.0), distributed under the Creative Commons Attribution 4.0 International licence, as recorded in the arXiv licence metadata for this exact version and shown on the abstract page https://arxiv.org/abs/2406.14414v4. Full rights notice: licenses/arxiv-2406.14414-CC-BY-4.0.txt.\par\smallskip

\noindent\textit{Editorial scope.} Counted unit: the Prop whose source label is \texttt{P:lower estimate} in the article (source0.tex lines 1278--1292, source label \texttt{P:lower estimate}), delivered together with the article's own complete original proof of it (source0.tex lines 1294--1395, one \texttt{proof} environment of 102 lines). No mathematics is written, repaired or re-derived: statement and proof are the article's own text line for line. Required context retained uncounted: none. Every definition, display and label of those lines is retained unchanged. Omitted from this package and not redistributed: 1--1277, 1293--1293, 1396--2307. Nothing inside a retained excerpt is omitted or altered: every retained body is delivered exactly as the article's lines are written, so no omission receipt of any kind accompanies this package. Citations: the retained excerpts carry no citation command at all, so no bibliography entry of the article is transcribed and no reference is redistributed. Reference adaptation: none was needed and none was made; every reference inside the delivered unit resolves inside it, so no unresolved reference or citation is produced. Printed slips kept literal: no printed word, symbol, hypothesis or number of the counted statement or of its proof was altered. Layout and class: the article's own class is \texttt{article} with packages including amsthm, amsmath, amsfonts, epsfig, amssymb, color, amscd, xy, pb-diagram, pb-xy, hyperref, tikz, eufrak; the wrapper is the lane's pinned \texttt{amsart} editorial wrapper at 10pt on A4, which restates the theorem-like environments the retained text needs and adds the article's own macro definitions that the retained text needs (\texttt{\textbackslash Ord}), with extra wrapper packages (bm, bbm, dsfont, xspace, cancel, mathtools). The delivered numbering is the unit's own. Assets: no retained span contains a figure environment, a graphics-inclusion command, a PDF-page inclusion, a table or a TikZ picture; no image file is copied into this package, the package declares no asset dependency and no image file is redistributed with it. Licence: version arXiv:2406.14414v4 of this article is distributed under the Creative Commons Attribution 4.0 International licence, as recorded in the arXiv licence metadata for this exact version and shown on the abstract page https://arxiv.org/abs/2406.14414v4. A full rights notice is delivered at licenses/arxiv-2406.14414-CC-BY-4.0.txt. Characters: every retained excerpt is pure ASCII, so the delivered engine, XeLaTeX with its default Latin Modern text font, reports no missing character.
\begin{Prop}
\label{P:lower estimate}
Suppose 
$$
0\neq H=\sum_{i=0}^{k-1}\sum_{m=0}^da_{m,i}A_i\Gamma_m\int^u.
$$ 	
is a HCP  from $Hcpc(k)$,  with $\Ord(H)=-u<0$. Then $H$ is totally free of $B_j$ iff the linear system of equations on $a_{m,i}$ holds:
\begin{equation}
\label{E:system_on_B_j}
\sum_{i=0}^{k-1}\sum_{m=0}^d\xi^{i(j-1)}(j-1)^ma_{m,i}=0
\end{equation}
for any $1\leq j\leq u$.

Moreover, if $H$ is totally free of $B_j$, then $\frac{u}{k}-1< Sdeg_A(H).$
\end{Prop}

\begin{proof}
    Suppose  $H$ is totally free of $B_j$. Since 
\begin{multline*}
	H\partial^u=\sum_{i=0}^{k-1}\sum_{m=0}^da_{m,i}A_i\Gamma_m(1-B_1-\cdots-B_{u})
	=\sum_{i=0}^{k-1}\sum_{m=0}^da_{m,i}A_i\Gamma_m-\sum_{j=1}^{u}\sum_{i=0}^{k-1}\sum_{m=0}^d\xi^{i(j-1)}(j-1)^ma_{m,i}B_j
\end{multline*}
is free of $B_j$, we get a linear equation system \eqref{E:system_on_B_j}. Vice versa, this system implies the total freeness of $B_j$: indeed, for any $l\ge 0$ $H\partial^l= H\partial^u D^{l-u}$, hence it is free of $B_j$ (for $l<0$  no terms with $B_j$ can appear). 

The second statement requires more technical arguments. 
   
 
\begin{lemma}
	\label{L:Tkdu full rank when u=(d+1)k}
	Suppose $T$ is the coefficient matrix of  system \eqref{E:system_on_B_j}. If $u=(d+1)k$, $T$ is a square matrix and we have $det(T)\neq 0$.
\end{lemma}

\begin{proof}
	Note that $T$ is in the shape of 
	$$
	\begin{pmatrix}
		1 & 1 & 1 & \cdots & 1 \\
		1 & \xi & \xi^2 & \cdots & \xi^{u-1} \\
		\cdots & \cdots & \cdots & \cdots & \cdots \\
		1 & \xi^{k-1} & \xi^{2(k-1)} & \cdots & \xi^{(k-1)(u-1)} \\
		\cdots & \cdots & \cdots & \cdots & \cdots \\
		0 & 1 & 2^d & \cdots & (u-1)^d \\
		0 & \xi & \xi^22^d & \cdots & \xi^{u-1}(u-1)^d \\
		\cdots & \cdots & \cdots & \cdots & \cdots \\
		0 & \xi^{k-1} & \xi^{2(k-1)}2^d & \cdots & \xi^{(k-1)(u-1)}(u-1)^d
	\end{pmatrix}
	$$
	suppose there exist $\overrightarrow{\alpha}=(\alpha_1,\dots,\alpha_{(d+1)k}),\alpha_l\in \tilde{K}, 1\leq l\leq k(d+1)$, such that 
	$$
	\overrightarrow{\alpha}T=0
	$$
	Then we have 
	\begin{multline*}
	(\alpha_{(d+1)k}\xi^{(k-1)(u-1)}+\cdots+\alpha_{(d+1)k-k+1})(u-1)^d+(\cdots)(u-1)^{d-1}
	+\cdots+(\alpha_k\xi^{(k-1)(u-1)}+\cdots+\alpha_1)=0
	\end{multline*}
	and 
	\begin{multline*}
		(\alpha_{(d+1)k}\xi^{(k-1)(u-2)}+\cdots+\alpha_{(d+1)k-k+1})(u-2)^d+(\cdots)(u-2)^{d-1}
		+\cdots+(\alpha_k\xi^{(k-1)(u-2)}+\cdots+\alpha_1)=0
	\end{multline*}
	\dots
\begin{multline*}
	(\alpha_{(d+1)k}\xi^{k-1}+\cdots+\alpha_{(d+1)k-k+1})+(\cdots)
	+\cdots+(\alpha_k\xi^{k-1}+\cdots+\alpha_1)=0
\end{multline*}

Now denote 
$$
\begin{cases}
\eta_0=\alpha_{(d+1)k}\xi^{(k-1)(u-1)}+\cdots+\alpha_{(d+1)k-k+1}
\\\eta_1=\alpha_{(d+1)k-k}\xi^{(k-1)(u-1)}+\cdots+\alpha_{(d+1)k-2k+1}
\\\cdots
\\\eta_d=\alpha_k\xi^{(k-1)(u-1)}+\cdots+\alpha_1
\end{cases}
$$
Notice that $\xi^k=1$, we have 
$$
\alpha_{(d+1)k}\xi^{(k-1)(u-k-1)}=\alpha_{(d+1)k}\xi^{(k-1)(u-1)}
$$
Hence we know 
$$
\begin{cases}
	(u-1)^d\eta_0+(u-1)^{d-1}+\cdots+\eta_d=0
	\\(u-k-1)^d\eta_0+(u-k-1)^{d-1}+\cdots+\eta_d=0
	\\\cdots
	\\(u-dk-1)^d\eta_0+(u-dk-1)^{d-2}\eta_1+\cdots+\eta_d=0
\end{cases}
$$
Considering $\eta_0,\eta_1,\dots,\eta_d$ as variables, we get a linear equation systems equipped with Vandermonde's coefficient matrix. Hence we know $\eta_0=\eta_1=\cdots=\eta_d=0$. In the same way, if we denote
$$
\begin{cases}
	\tilde{\eta}_0=\alpha_{(d+1)k}\xi^{(k-1)(u-2)}+\cdots+\alpha_{(d+1)k-k+1}
	\\\tilde{\eta}_1=\alpha_{(d+1)k-k}\xi^{(k-1)(u-2)}+\cdots+\alpha_{(d+1)k-2k+1}
	\\\cdots
	\\\tilde{\eta}_d=\alpha_k\xi^{(k-1)(u-2)}+\cdots+\alpha_1
\end{cases}
$$
we will also get $\tilde{\eta}_0=\tilde{\eta}_1=\cdots=\tilde{\eta}_d=0$. We play with this game for $k$ times, there will be $k$ numbers of equation for $(\alpha_{1+mk},\alpha_{2+mk},\dots,\alpha_{k+mk})$ for any $0\leq m\leq d$, in the form of 
$$
(\alpha_{1+mk},\alpha_{2+mk},\dots,\alpha_{k+mk})\begin{pmatrix}
	1 & 1 & \cdots & 1 \\
	1 & \xi^1 & \cdots & \xi^{k-1} \\
	\cdots & \cdots & \cdots & \cdots \\
	1 & \xi^{k-1} & \cdots & \xi^{(k-1)(k-1)}
\end{pmatrix}=0
$$
The coefficient matrix is again in Vandermonde's shape, so there are only $0$ solutions, this means $\overrightarrow{\alpha}=0$. Hence $det(T)\neq 0$.
\end{proof}
   
    
    
	Summing up, $T$ induces a linear map: 
	$$
	T:\tilde{K}^{k(d+1)} \rightarrow \tilde{K}^{u}.
	$$
If $d=Sdeg_A(H)\leq  \frac{u}{k}-1$, i.e. $(d+1)k\leq  u$, then  $T$ is injective (as its matrix consists of linearly independent raws by lemma \ref{L:Tkdu full rank when u=(d+1)k}), this means the linear equation system has  only zero solution, this contradicts to $H\neq 0$.
\end{proof}

\end{document}
